% App H: General PGT+EM Lagrangian — sector-by-sector enumeration.
% Forward-reference target from §2.1 for the explicit invariant
% lists and parameter counts. The theory-background formalism
% (vierbein, spin connection, irreducible torsion, ghost-restriction
% concepts) lives in App F (pgt_background.tex).

\section{General PGT+EM Lagrangian: sector-by-sector enumeration}
\label{PGTEnumeration}

\paragraph*{Linear sector}
\label{LinearSector}
% TODO ~80 w. Three independent linear invariants: Einstein--Hilbert
% term R/\kappa^2, cosmological constant \Lambda, and the Holst term
% \varepsilon^{\mu\nu\rho\sigma} \tilde{R}_{\mu\nu\rho\sigma} with
% Barbero--Immirzi parameter \gamma_BI. Cite holst1995barbero,
% kibble1961lorentz.

\paragraph*{Parity-even quadratic sector}
\label{ParityEvenQuadratic}
% TODO ~250 w + 1 equation block. Eleven independent invariants:
% six \tilde{R}^2 (\alpha_1 ... \alpha_6 in research-doc convention,
% \beta_1 ... \beta_6 in §2.1 manuscript convention), three
% \mathcal{T}^2 (three irreducible torsion-squared, written in
% direct basis {I_1, I_2, I_3} as in §2.1), one F^2 (Maxwell), and
% one \tilde{R}_{[\mu\nu]} F^{\mu\nu} cross-term (\delta_1).
% Three of the six \tilde{R}^2 invariants reduce in the torsion-free
% case via Bianchi + pair-exchange symmetry, with Gauss--Bonnet
% further reducing to two; with torsion all six are independent
% (citing the project's research/lagrangian_enumeration/
% general_quadratic_lagrangian.tex enumeration). Cite lin2019pgt,
% blagojevic2018hamiltonian, blagojevic2003lectures,
% barker2024poincare, sezgin1980ghostfree, aoki2020nonlinear.

\paragraph*{Parity-odd quadratic sector}
\label{ParityOddQuadratic}
% TODO ~150 w. Twenty additional physical parity-odd quadratic
% invariants from \varepsilon-weighted contractions of curvature,
% torsion, and EM field strength: e.g.\ \varepsilon^{\mu\nu\rho\sigma}
% T_{\mu\nu\alpha} T_{\rho\sigma}{}^\alpha, \tilde{R}_{[\mu\nu]}
% \tilde{F}^{\mu\nu} (Hodge-dual EM cross), Holst-like
% \tilde{R}_{\mu\nu\rho\sigma}\tilde{R}^{\rho\sigma\mu\nu}.
% Karananas 2014 + Blagojevic--Cvetkovic 2018 give the parity-odd
% mass-spectrum analyses. Cite karananas2014parity,
% blagojevic2018hamiltonian.

\paragraph*{Ghost-free derivative extensions}
\label{DerivativeExtensions}
% TODO ~100 w. Ghost-free higher-derivative completions add at
% most ~29 parity-even and ~80 parity-odd invariants (upper bounds
% before Bianchi reduction), as enumerated in
% research/lagrangian_enumeration/general_quadratic_lagrangian.tex.
% Aoki--Yokoyama 2020 provides the nonlinearly-ghost-free
% higher-curvature analysis. Cite aoki2020nonlinear.

\paragraph*{Cubic completions}
\label{CubicCompletions}
% TODO ~80 w. Approximately 65 safe cubic operators relevant for
% (i) ghost elimination in propagating-torsion theories and (ii)
% linearisation around magnetic backgrounds (in which the cubic
% terms contribute at quadratic order in the perturbation).
% Bahamonde et al. 2025 constructs ghost-free cubic PGT
% cosmologies. Cite bahamonde2025cubic.

\paragraph*{Chern--Simons torsion--EM couplings}
\label{ChernSimonsCouplings}
% TODO ~120 w. Parity-odd couplings involving the gauge potential
% A_\mu directly (rather than only F_{\mu\nu}): the trace-torsion
% Chern--Simons \chi^{CS}_2 \varepsilon^{\mu\nu\rho\sigma} T_\mu A_\nu
% \partial_\rho A_\sigma; the full-torsion Chern--Simons
% \chi^{CS}_3 \varepsilon^{\mu\nu\rho\sigma} T_{\mu\nu}{}^\alpha
% A_\alpha \partial_\rho A_\sigma (one term + six tensor variants).
% Hehl--Obukhov 2000 gave the original axion/birefringence
% framework; Itin--Hehl 2003 the systematic T^2 F^2 classification;
% Rubilar 2003 the torsion-birefringence treatment; Obukhov et
% al.\ 2024 the recent classical-electrodynamics + spin-torsion
% analysis. Cite hehl2000electromagnetic, itin2004maxwell,
% rubilar2003torsion, obukhov2024electrodynamics.

\paragraph*{Direct--irreducible basis translation}
\label{DirectIrreducibleTranslation}
% TODO ~150 w + 1 equation block. The 3x3 invertible
% transformation between the direct basis {I_1, I_2, I_3} of
% torsion-squared scalar contractions and the irreducible basis
% {q^2, T^2, S^2}; the explicit conversion factors. Convention
% comparison table: research/lagrangian_enumeration uses (\alpha
% for \tilde{R}^2, \beta for T^2); the §2.1 manuscript convention
% uses (\alpha for T^2, \beta for \tilde{R}^2) matching TIDAL's
% code labels (alpha1, alpha2, alpha3 in TOML files). Cite
% shapiro2002torsion §2.3.
\begin{align}
  % TODO: 3x3 transformation matrix between {I_1, I_2, I_3}
  % and {q^2, T^2, S^2}
  \label{DirectIrreducibleTransform}
\end{align}
